\documentclass[addpoints,spanish, 12pt,a4paper]{exam}
%\documentclass[answers, spanish, 12pt,a4paper]{exam}
\printanswers
\pointpoints{punto}{puntos}
\hpword{Puntos:}
\vpword{Puntos:}
\htword{Total}
\vtword{Total}
\hsword{Resultado:}
\hqword{Ejercicio:}
\vqword{Ejercicio:}

\usepackage[utf8]{inputenc}
\usepackage[spanish]{babel}
\usepackage{eurosym}
%\usepackage[spanish,es-lcroman, es-tabla, es-noshorthands]{babel}
        \usepackage{pgf}


\usepackage[margin=1in]{geometry}
\usepackage{amsmath,amssymb}
\usepackage{multicol}
\usepackage{yhmath}
\usepackage{pdflscape}

\pointsinrightmargin % Para poner las puntuaciones a la derecha. Se puede cambiar. Si se comenta, sale a la izquierda.
\extrawidth{-2.4cm} %Un poquito más de margen por si ponemos textos largos.
\marginpointname{ \emph{\points}}

\usepackage{graphicx}
\graphicspath{{../img/}}

\newcommand{\class}{1º Bachillerato CCSS}
\newcommand{\examdate}{\today}
\newcommand{\examnum}{Recuperación 2ª Evaluación}
\newcommand{\tipo}{B}


\newcommand{\timelimit}{50 minutos}

\renewcommand{\solutiontitle}{\noindent\textbf{Solución:}\enspace}


\pagestyle{head}
\firstpageheader{\includegraphics[width=0.2\columnwidth]{header_left}}{\textbf{Departamento de Matemáticas\linebreak \class}\linebreak \examnum}{\includegraphics[width=0.1\columnwidth]{header_right}}
\runningheader{\class}{\examnum}{Página \thepage\ of \numpages}
\runningheadrule


\begin{document}

\noindent
\begin{tabular*}{\textwidth}{l @{\extracolsep{\fill}} r @{\extracolsep{6pt}} }
\textbf{Nombre:} \makebox[3.5in]{\hrulefill} & \textbf{Fecha:}\makebox[1in]{\hrulefill} \\
 & \\
\textbf{Tiempo: \timelimit} & Tipo: \tipo
\end{tabular*}
\rule[2ex]{\textwidth}{2pt}
Esta prueba tiene \numquestions\ ejercicios. La puntuación máxima es de \numpoints.
La nota final de la prueba será la parte proporcional de la puntuación obtenida sobre la puntuación máxima.

\begin{center}


\addpoints
 %\gradetable[h][questions]
	\pointtable[h][questions]
\end{center}

\noindent
\rule[2ex]{\textwidth}{2pt}

\begin{questions}


        \question au31e03 - La temperatura media en los meses de invierno en varias ciudades y el gasto medio por habitante en
calefacción ha sido\\\begin{tabular}{lrrrr}
\hline
                      &   0 &   1 &   2 &   3 \\
\hline
 Temperatura (grados) &  10 &  14 &  17 &  20 \\
 Gasto (euros)        & 150 & 102 &  55 &  18 \\
\hline
\end{tabular}
        \begin{parts} \part[1] Haz una tabla de frecuencias con los datos que necesites para hace el resto de apartados  \begin{solution}   \begin{tabular}{rrrrrr}
\hline
    &   x &   y &   xy &   x2 &    y2 \\
\hline
  0 &  10 & 150 & 1500 &  100 & 22500 \\
  1 &  14 & 102 & 1428 &  196 & 10404 \\
  2 &  17 &  55 &  935 &  289 &  3025 \\
  3 &  20 &  18 &  360 &  400 &   324 \\
  4 &  61 & 325 & 4223 &  985 & 36253 \\
\hline
\end{tabular}   \end{solution} \part[1] Calcula el gasto medio  \begin{solution}   {'media': 81.25}   \end{solution} \part[2] Halla el coeficiente de correlación lineal e interprétalo  \begin{solution}   {'media de x': 15.25, 'desviación de x': 3.6996621467371855, 'media de y': 81.25, 'desviación de y': 49.61539579606314, 'covarianza': -183.3125, 'coeficiente de correlación': -0.9986505695692516}   \end{solution} \part[1] Estima el gasto medio por habitante de una ciudad si la temperatura media hubiera sido 8ºC  \begin{solution}   $y = - 13.3926940639269 x + 285.488584474886$ \\\resizebox {0.3\textwidth}{!}{%% Creator: Matplotlib, PGF backend
%%
%% To include the figure in your LaTeX document, write
%%   \input{<filename>.pgf}
%%
%% Make sure the required packages are loaded in your preamble
%%   \usepackage{pgf}
%%
%% Figures using additional raster images can only be included by \input if
%% they are in the same directory as the main LaTeX file. For loading figures
%% from other directories you can use the `import` package
%%   \usepackage{import}
%% and then include the figures with
%%   \import{<path to file>}{<filename>.pgf}
%%
%% Matplotlib used the following preamble
%%   \usepackage{fontspec}
%%   \setmainfont{DejaVuSerif.ttf}[Path=/home/hp/Mis_aplicaciones/anaconda3/lib/python3.6/site-packages/matplotlib/mpl-data/fonts/ttf/]
%%   \setsansfont{DejaVuSans.ttf}[Path=/home/hp/Mis_aplicaciones/anaconda3/lib/python3.6/site-packages/matplotlib/mpl-data/fonts/ttf/]
%%   \setmonofont{DejaVuSansMono.ttf}[Path=/home/hp/Mis_aplicaciones/anaconda3/lib/python3.6/site-packages/matplotlib/mpl-data/fonts/ttf/]
%%
\begingroup%
\makeatletter%
\begin{pgfpicture}%
\pgfpathrectangle{\pgfpointorigin}{\pgfqpoint{6.000000in}{4.000000in}}%
\pgfusepath{use as bounding box, clip}%
\begin{pgfscope}%
\pgfsetbuttcap%
\pgfsetmiterjoin%
\definecolor{currentfill}{rgb}{1.000000,1.000000,1.000000}%
\pgfsetfillcolor{currentfill}%
\pgfsetlinewidth{0.000000pt}%
\definecolor{currentstroke}{rgb}{1.000000,1.000000,1.000000}%
\pgfsetstrokecolor{currentstroke}%
\pgfsetdash{}{0pt}%
\pgfpathmoveto{\pgfqpoint{0.000000in}{0.000000in}}%
\pgfpathlineto{\pgfqpoint{6.000000in}{0.000000in}}%
\pgfpathlineto{\pgfqpoint{6.000000in}{4.000000in}}%
\pgfpathlineto{\pgfqpoint{0.000000in}{4.000000in}}%
\pgfpathclose%
\pgfusepath{fill}%
\end{pgfscope}%
\begin{pgfscope}%
\pgfsetbuttcap%
\pgfsetmiterjoin%
\definecolor{currentfill}{rgb}{1.000000,1.000000,1.000000}%
\pgfsetfillcolor{currentfill}%
\pgfsetlinewidth{0.000000pt}%
\definecolor{currentstroke}{rgb}{0.000000,0.000000,0.000000}%
\pgfsetstrokecolor{currentstroke}%
\pgfsetstrokeopacity{0.000000}%
\pgfsetdash{}{0pt}%
\pgfpathmoveto{\pgfqpoint{0.750000in}{0.500000in}}%
\pgfpathlineto{\pgfqpoint{5.400000in}{0.500000in}}%
\pgfpathlineto{\pgfqpoint{5.400000in}{3.520000in}}%
\pgfpathlineto{\pgfqpoint{0.750000in}{3.520000in}}%
\pgfpathclose%
\pgfusepath{fill}%
\end{pgfscope}%
\begin{pgfscope}%
\pgfpathrectangle{\pgfqpoint{0.750000in}{0.500000in}}{\pgfqpoint{4.650000in}{3.020000in}}%
\pgfusepath{clip}%
\pgfsetbuttcap%
\pgfsetroundjoin%
\definecolor{currentfill}{rgb}{0.121569,0.466667,0.705882}%
\pgfsetfillcolor{currentfill}%
\pgfsetfillopacity{0.800000}%
\pgfsetlinewidth{1.003750pt}%
\definecolor{currentstroke}{rgb}{0.121569,0.466667,0.705882}%
\pgfsetstrokecolor{currentstroke}%
\pgfsetstrokeopacity{0.800000}%
\pgfsetdash{}{0pt}%
\pgfpathmoveto{\pgfqpoint{0.967699in}{2.979571in}}%
\pgfpathcurveto{\pgfqpoint{0.978749in}{2.979571in}}{\pgfqpoint{0.989348in}{2.983962in}}{\pgfqpoint{0.997162in}{2.991775in}}%
\pgfpathcurveto{\pgfqpoint{1.004975in}{2.999589in}}{\pgfqpoint{1.009365in}{3.010188in}}{\pgfqpoint{1.009365in}{3.021238in}}%
\pgfpathcurveto{\pgfqpoint{1.009365in}{3.032288in}}{\pgfqpoint{1.004975in}{3.042887in}}{\pgfqpoint{0.997162in}{3.050701in}}%
\pgfpathcurveto{\pgfqpoint{0.989348in}{3.058515in}}{\pgfqpoint{0.978749in}{3.062905in}}{\pgfqpoint{0.967699in}{3.062905in}}%
\pgfpathcurveto{\pgfqpoint{0.956649in}{3.062905in}}{\pgfqpoint{0.946050in}{3.058515in}}{\pgfqpoint{0.938236in}{3.050701in}}%
\pgfpathcurveto{\pgfqpoint{0.930422in}{3.042887in}}{\pgfqpoint{0.926032in}{3.032288in}}{\pgfqpoint{0.926032in}{3.021238in}}%
\pgfpathcurveto{\pgfqpoint{0.926032in}{3.010188in}}{\pgfqpoint{0.930422in}{2.999589in}}{\pgfqpoint{0.938236in}{2.991775in}}%
\pgfpathcurveto{\pgfqpoint{0.946050in}{2.983962in}}{\pgfqpoint{0.956649in}{2.979571in}}{\pgfqpoint{0.967699in}{2.979571in}}%
\pgfpathclose%
\pgfusepath{stroke,fill}%
\end{pgfscope}%
\begin{pgfscope}%
\pgfpathrectangle{\pgfqpoint{0.750000in}{0.500000in}}{\pgfqpoint{4.650000in}{3.020000in}}%
\pgfusepath{clip}%
\pgfsetbuttcap%
\pgfsetroundjoin%
\definecolor{currentfill}{rgb}{0.121569,0.466667,0.705882}%
\pgfsetfillcolor{currentfill}%
\pgfsetfillopacity{0.800000}%
\pgfsetlinewidth{1.003750pt}%
\definecolor{currentstroke}{rgb}{0.121569,0.466667,0.705882}%
\pgfsetstrokecolor{currentstroke}%
\pgfsetstrokeopacity{0.800000}%
\pgfsetdash{}{0pt}%
\pgfpathmoveto{\pgfqpoint{2.653540in}{2.217172in}}%
\pgfpathcurveto{\pgfqpoint{2.664590in}{2.217172in}}{\pgfqpoint{2.675189in}{2.221562in}}{\pgfqpoint{2.683003in}{2.229376in}}%
\pgfpathcurveto{\pgfqpoint{2.690816in}{2.237190in}}{\pgfqpoint{2.695206in}{2.247789in}}{\pgfqpoint{2.695206in}{2.258839in}}%
\pgfpathcurveto{\pgfqpoint{2.695206in}{2.269889in}}{\pgfqpoint{2.690816in}{2.280488in}}{\pgfqpoint{2.683003in}{2.288301in}}%
\pgfpathcurveto{\pgfqpoint{2.675189in}{2.296115in}}{\pgfqpoint{2.664590in}{2.300505in}}{\pgfqpoint{2.653540in}{2.300505in}}%
\pgfpathcurveto{\pgfqpoint{2.642490in}{2.300505in}}{\pgfqpoint{2.631891in}{2.296115in}}{\pgfqpoint{2.624077in}{2.288301in}}%
\pgfpathcurveto{\pgfqpoint{2.616263in}{2.280488in}}{\pgfqpoint{2.611873in}{2.269889in}}{\pgfqpoint{2.611873in}{2.258839in}}%
\pgfpathcurveto{\pgfqpoint{2.611873in}{2.247789in}}{\pgfqpoint{2.616263in}{2.237190in}}{\pgfqpoint{2.624077in}{2.229376in}}%
\pgfpathcurveto{\pgfqpoint{2.631891in}{2.221562in}}{\pgfqpoint{2.642490in}{2.217172in}}{\pgfqpoint{2.653540in}{2.217172in}}%
\pgfpathclose%
\pgfusepath{stroke,fill}%
\end{pgfscope}%
\begin{pgfscope}%
\pgfpathrectangle{\pgfqpoint{0.750000in}{0.500000in}}{\pgfqpoint{4.650000in}{3.020000in}}%
\pgfusepath{clip}%
\pgfsetbuttcap%
\pgfsetroundjoin%
\definecolor{currentfill}{rgb}{0.121569,0.466667,0.705882}%
\pgfsetfillcolor{currentfill}%
\pgfsetfillopacity{0.800000}%
\pgfsetlinewidth{1.003750pt}%
\definecolor{currentstroke}{rgb}{0.121569,0.466667,0.705882}%
\pgfsetstrokecolor{currentstroke}%
\pgfsetstrokeopacity{0.800000}%
\pgfsetdash{}{0pt}%
\pgfpathmoveto{\pgfqpoint{3.917920in}{1.470656in}}%
\pgfpathcurveto{\pgfqpoint{3.928971in}{1.470656in}}{\pgfqpoint{3.939570in}{1.475046in}}{\pgfqpoint{3.947383in}{1.482860in}}%
\pgfpathcurveto{\pgfqpoint{3.955197in}{1.490673in}}{\pgfqpoint{3.959587in}{1.501272in}}{\pgfqpoint{3.959587in}{1.512323in}}%
\pgfpathcurveto{\pgfqpoint{3.959587in}{1.523373in}}{\pgfqpoint{3.955197in}{1.533972in}}{\pgfqpoint{3.947383in}{1.541785in}}%
\pgfpathcurveto{\pgfqpoint{3.939570in}{1.549599in}}{\pgfqpoint{3.928971in}{1.553989in}}{\pgfqpoint{3.917920in}{1.553989in}}%
\pgfpathcurveto{\pgfqpoint{3.906870in}{1.553989in}}{\pgfqpoint{3.896271in}{1.549599in}}{\pgfqpoint{3.888458in}{1.541785in}}%
\pgfpathcurveto{\pgfqpoint{3.880644in}{1.533972in}}{\pgfqpoint{3.876254in}{1.523373in}}{\pgfqpoint{3.876254in}{1.512323in}}%
\pgfpathcurveto{\pgfqpoint{3.876254in}{1.501272in}}{\pgfqpoint{3.880644in}{1.490673in}}{\pgfqpoint{3.888458in}{1.482860in}}%
\pgfpathcurveto{\pgfqpoint{3.896271in}{1.475046in}}{\pgfqpoint{3.906870in}{1.470656in}}{\pgfqpoint{3.917920in}{1.470656in}}%
\pgfpathclose%
\pgfusepath{stroke,fill}%
\end{pgfscope}%
\begin{pgfscope}%
\pgfpathrectangle{\pgfqpoint{0.750000in}{0.500000in}}{\pgfqpoint{4.650000in}{3.020000in}}%
\pgfusepath{clip}%
\pgfsetbuttcap%
\pgfsetroundjoin%
\definecolor{currentfill}{rgb}{0.121569,0.466667,0.705882}%
\pgfsetfillcolor{currentfill}%
\pgfsetfillopacity{0.800000}%
\pgfsetlinewidth{1.003750pt}%
\definecolor{currentstroke}{rgb}{0.121569,0.466667,0.705882}%
\pgfsetstrokecolor{currentstroke}%
\pgfsetstrokeopacity{0.800000}%
\pgfsetdash{}{0pt}%
\pgfpathmoveto{\pgfqpoint{5.182301in}{0.882973in}}%
\pgfpathcurveto{\pgfqpoint{5.193351in}{0.882973in}}{\pgfqpoint{5.203950in}{0.887363in}}{\pgfqpoint{5.211764in}{0.895177in}}%
\pgfpathcurveto{\pgfqpoint{5.219578in}{0.902991in}}{\pgfqpoint{5.223968in}{0.913590in}}{\pgfqpoint{5.223968in}{0.924640in}}%
\pgfpathcurveto{\pgfqpoint{5.223968in}{0.935690in}}{\pgfqpoint{5.219578in}{0.946289in}}{\pgfqpoint{5.211764in}{0.954102in}}%
\pgfpathcurveto{\pgfqpoint{5.203950in}{0.961916in}}{\pgfqpoint{5.193351in}{0.966306in}}{\pgfqpoint{5.182301in}{0.966306in}}%
\pgfpathcurveto{\pgfqpoint{5.171251in}{0.966306in}}{\pgfqpoint{5.160652in}{0.961916in}}{\pgfqpoint{5.152838in}{0.954102in}}%
\pgfpathcurveto{\pgfqpoint{5.145025in}{0.946289in}}{\pgfqpoint{5.140635in}{0.935690in}}{\pgfqpoint{5.140635in}{0.924640in}}%
\pgfpathcurveto{\pgfqpoint{5.140635in}{0.913590in}}{\pgfqpoint{5.145025in}{0.902991in}}{\pgfqpoint{5.152838in}{0.895177in}}%
\pgfpathcurveto{\pgfqpoint{5.160652in}{0.887363in}}{\pgfqpoint{5.171251in}{0.882973in}}{\pgfqpoint{5.182301in}{0.882973in}}%
\pgfpathclose%
\pgfusepath{stroke,fill}%
\end{pgfscope}%
\begin{pgfscope}%
\pgfpathrectangle{\pgfqpoint{0.750000in}{0.500000in}}{\pgfqpoint{4.650000in}{3.020000in}}%
\pgfusepath{clip}%
\pgfsetbuttcap%
\pgfsetroundjoin%
\definecolor{currentfill}{rgb}{0.121569,0.466667,0.705882}%
\pgfsetfillcolor{currentfill}%
\pgfsetfillopacity{0.150000}%
\pgfsetlinewidth{0.000000pt}%
\definecolor{currentstroke}{rgb}{0.000000,0.000000,0.000000}%
\pgfsetstrokecolor{currentstroke}%
\pgfsetdash{}{0pt}%
\pgfpathmoveto{\pgfqpoint{0.750000in}{3.382727in}}%
\pgfpathlineto{\pgfqpoint{0.750000in}{2.984769in}}%
\pgfpathlineto{\pgfqpoint{0.796970in}{2.962937in}}%
\pgfpathlineto{\pgfqpoint{0.843939in}{2.941106in}}%
\pgfpathlineto{\pgfqpoint{0.890909in}{2.919274in}}%
\pgfpathlineto{\pgfqpoint{0.937879in}{2.897443in}}%
\pgfpathlineto{\pgfqpoint{0.984848in}{2.875611in}}%
\pgfpathlineto{\pgfqpoint{1.031818in}{2.853780in}}%
\pgfpathlineto{\pgfqpoint{1.078788in}{2.831949in}}%
\pgfpathlineto{\pgfqpoint{1.125758in}{2.810117in}}%
\pgfpathlineto{\pgfqpoint{1.172727in}{2.788286in}}%
\pgfpathlineto{\pgfqpoint{1.219697in}{2.766454in}}%
\pgfpathlineto{\pgfqpoint{1.266667in}{2.744623in}}%
\pgfpathlineto{\pgfqpoint{1.313636in}{2.722791in}}%
\pgfpathlineto{\pgfqpoint{1.360606in}{2.700960in}}%
\pgfpathlineto{\pgfqpoint{1.407576in}{2.679128in}}%
\pgfpathlineto{\pgfqpoint{1.454545in}{2.657297in}}%
\pgfpathlineto{\pgfqpoint{1.501515in}{2.635465in}}%
\pgfpathlineto{\pgfqpoint{1.548485in}{2.613634in}}%
\pgfpathlineto{\pgfqpoint{1.595455in}{2.591802in}}%
\pgfpathlineto{\pgfqpoint{1.642424in}{2.569971in}}%
\pgfpathlineto{\pgfqpoint{1.689394in}{2.548139in}}%
\pgfpathlineto{\pgfqpoint{1.736364in}{2.526308in}}%
\pgfpathlineto{\pgfqpoint{1.783333in}{2.504477in}}%
\pgfpathlineto{\pgfqpoint{1.830303in}{2.482645in}}%
\pgfpathlineto{\pgfqpoint{1.877273in}{2.460814in}}%
\pgfpathlineto{\pgfqpoint{1.924242in}{2.438982in}}%
\pgfpathlineto{\pgfqpoint{1.971212in}{2.417151in}}%
\pgfpathlineto{\pgfqpoint{2.018182in}{2.395319in}}%
\pgfpathlineto{\pgfqpoint{2.065152in}{2.373488in}}%
\pgfpathlineto{\pgfqpoint{2.112121in}{2.351656in}}%
\pgfpathlineto{\pgfqpoint{2.159091in}{2.329825in}}%
\pgfpathlineto{\pgfqpoint{2.206061in}{2.307993in}}%
\pgfpathlineto{\pgfqpoint{2.253030in}{2.286162in}}%
\pgfpathlineto{\pgfqpoint{2.300000in}{2.264330in}}%
\pgfpathlineto{\pgfqpoint{2.346970in}{2.242499in}}%
\pgfpathlineto{\pgfqpoint{2.393939in}{2.220667in}}%
\pgfpathlineto{\pgfqpoint{2.440909in}{2.198836in}}%
\pgfpathlineto{\pgfqpoint{2.487879in}{2.177005in}}%
\pgfpathlineto{\pgfqpoint{2.534848in}{2.155173in}}%
\pgfpathlineto{\pgfqpoint{2.581818in}{2.133342in}}%
\pgfpathlineto{\pgfqpoint{2.628788in}{2.111510in}}%
\pgfpathlineto{\pgfqpoint{2.675758in}{2.089679in}}%
\pgfpathlineto{\pgfqpoint{2.722727in}{2.067847in}}%
\pgfpathlineto{\pgfqpoint{2.769697in}{2.046016in}}%
\pgfpathlineto{\pgfqpoint{2.816667in}{2.024184in}}%
\pgfpathlineto{\pgfqpoint{2.863636in}{2.002353in}}%
\pgfpathlineto{\pgfqpoint{2.910606in}{1.980521in}}%
\pgfpathlineto{\pgfqpoint{2.957576in}{1.958690in}}%
\pgfpathlineto{\pgfqpoint{3.004545in}{1.936858in}}%
\pgfpathlineto{\pgfqpoint{3.051515in}{1.915027in}}%
\pgfpathlineto{\pgfqpoint{3.098485in}{1.893195in}}%
\pgfpathlineto{\pgfqpoint{3.145455in}{1.871364in}}%
\pgfpathlineto{\pgfqpoint{3.192424in}{1.849533in}}%
\pgfpathlineto{\pgfqpoint{3.239394in}{1.827701in}}%
\pgfpathlineto{\pgfqpoint{3.286364in}{1.805870in}}%
\pgfpathlineto{\pgfqpoint{3.333333in}{1.784038in}}%
\pgfpathlineto{\pgfqpoint{3.380303in}{1.762207in}}%
\pgfpathlineto{\pgfqpoint{3.427273in}{1.740375in}}%
\pgfpathlineto{\pgfqpoint{3.474242in}{1.718544in}}%
\pgfpathlineto{\pgfqpoint{3.521212in}{1.696712in}}%
\pgfpathlineto{\pgfqpoint{3.568182in}{1.674881in}}%
\pgfpathlineto{\pgfqpoint{3.615152in}{1.653049in}}%
\pgfpathlineto{\pgfqpoint{3.662121in}{1.631218in}}%
\pgfpathlineto{\pgfqpoint{3.709091in}{1.609386in}}%
\pgfpathlineto{\pgfqpoint{3.756061in}{1.587555in}}%
\pgfpathlineto{\pgfqpoint{3.803030in}{1.565723in}}%
\pgfpathlineto{\pgfqpoint{3.850000in}{1.543892in}}%
\pgfpathlineto{\pgfqpoint{3.896970in}{1.522061in}}%
\pgfpathlineto{\pgfqpoint{3.943939in}{1.496961in}}%
\pgfpathlineto{\pgfqpoint{3.990909in}{1.469229in}}%
\pgfpathlineto{\pgfqpoint{4.037879in}{1.441497in}}%
\pgfpathlineto{\pgfqpoint{4.084848in}{1.413765in}}%
\pgfpathlineto{\pgfqpoint{4.131818in}{1.386033in}}%
\pgfpathlineto{\pgfqpoint{4.178788in}{1.358301in}}%
\pgfpathlineto{\pgfqpoint{4.225758in}{1.330569in}}%
\pgfpathlineto{\pgfqpoint{4.272727in}{1.302837in}}%
\pgfpathlineto{\pgfqpoint{4.319697in}{1.275106in}}%
\pgfpathlineto{\pgfqpoint{4.366667in}{1.247374in}}%
\pgfpathlineto{\pgfqpoint{4.413636in}{1.219642in}}%
\pgfpathlineto{\pgfqpoint{4.460606in}{1.191910in}}%
\pgfpathlineto{\pgfqpoint{4.507576in}{1.164178in}}%
\pgfpathlineto{\pgfqpoint{4.554545in}{1.136446in}}%
\pgfpathlineto{\pgfqpoint{4.601515in}{1.108714in}}%
\pgfpathlineto{\pgfqpoint{4.648485in}{1.080983in}}%
\pgfpathlineto{\pgfqpoint{4.695455in}{1.053251in}}%
\pgfpathlineto{\pgfqpoint{4.742424in}{1.025519in}}%
\pgfpathlineto{\pgfqpoint{4.789394in}{0.997787in}}%
\pgfpathlineto{\pgfqpoint{4.836364in}{0.970055in}}%
\pgfpathlineto{\pgfqpoint{4.883333in}{0.942323in}}%
\pgfpathlineto{\pgfqpoint{4.930303in}{0.914591in}}%
\pgfpathlineto{\pgfqpoint{4.977273in}{0.886860in}}%
\pgfpathlineto{\pgfqpoint{5.024242in}{0.859128in}}%
\pgfpathlineto{\pgfqpoint{5.071212in}{0.831396in}}%
\pgfpathlineto{\pgfqpoint{5.118182in}{0.803664in}}%
\pgfpathlineto{\pgfqpoint{5.165152in}{0.775932in}}%
\pgfpathlineto{\pgfqpoint{5.212121in}{0.748200in}}%
\pgfpathlineto{\pgfqpoint{5.259091in}{0.720468in}}%
\pgfpathlineto{\pgfqpoint{5.306061in}{0.692736in}}%
\pgfpathlineto{\pgfqpoint{5.353030in}{0.665005in}}%
\pgfpathlineto{\pgfqpoint{5.400000in}{0.637273in}}%
\pgfpathlineto{\pgfqpoint{5.400000in}{1.016788in}}%
\pgfpathlineto{\pgfqpoint{5.400000in}{1.016788in}}%
\pgfpathlineto{\pgfqpoint{5.353030in}{1.038030in}}%
\pgfpathlineto{\pgfqpoint{5.306061in}{1.059271in}}%
\pgfpathlineto{\pgfqpoint{5.259091in}{1.080512in}}%
\pgfpathlineto{\pgfqpoint{5.212121in}{1.101754in}}%
\pgfpathlineto{\pgfqpoint{5.165152in}{1.122995in}}%
\pgfpathlineto{\pgfqpoint{5.118182in}{1.144237in}}%
\pgfpathlineto{\pgfqpoint{5.071212in}{1.165478in}}%
\pgfpathlineto{\pgfqpoint{5.024242in}{1.186720in}}%
\pgfpathlineto{\pgfqpoint{4.977273in}{1.207961in}}%
\pgfpathlineto{\pgfqpoint{4.930303in}{1.229202in}}%
\pgfpathlineto{\pgfqpoint{4.883333in}{1.250444in}}%
\pgfpathlineto{\pgfqpoint{4.836364in}{1.271685in}}%
\pgfpathlineto{\pgfqpoint{4.789394in}{1.292927in}}%
\pgfpathlineto{\pgfqpoint{4.742424in}{1.314168in}}%
\pgfpathlineto{\pgfqpoint{4.695455in}{1.335410in}}%
\pgfpathlineto{\pgfqpoint{4.648485in}{1.356651in}}%
\pgfpathlineto{\pgfqpoint{4.601515in}{1.377892in}}%
\pgfpathlineto{\pgfqpoint{4.554545in}{1.399134in}}%
\pgfpathlineto{\pgfqpoint{4.507576in}{1.420375in}}%
\pgfpathlineto{\pgfqpoint{4.460606in}{1.441617in}}%
\pgfpathlineto{\pgfqpoint{4.413636in}{1.462858in}}%
\pgfpathlineto{\pgfqpoint{4.366667in}{1.484100in}}%
\pgfpathlineto{\pgfqpoint{4.319697in}{1.505341in}}%
\pgfpathlineto{\pgfqpoint{4.272727in}{1.526582in}}%
\pgfpathlineto{\pgfqpoint{4.225758in}{1.547824in}}%
\pgfpathlineto{\pgfqpoint{4.178788in}{1.569065in}}%
\pgfpathlineto{\pgfqpoint{4.131818in}{1.590307in}}%
\pgfpathlineto{\pgfqpoint{4.084848in}{1.611548in}}%
\pgfpathlineto{\pgfqpoint{4.037879in}{1.632790in}}%
\pgfpathlineto{\pgfqpoint{3.990909in}{1.654031in}}%
\pgfpathlineto{\pgfqpoint{3.943939in}{1.675272in}}%
\pgfpathlineto{\pgfqpoint{3.896970in}{1.696514in}}%
\pgfpathlineto{\pgfqpoint{3.850000in}{1.717755in}}%
\pgfpathlineto{\pgfqpoint{3.803030in}{1.738997in}}%
\pgfpathlineto{\pgfqpoint{3.756061in}{1.760238in}}%
\pgfpathlineto{\pgfqpoint{3.709091in}{1.781480in}}%
\pgfpathlineto{\pgfqpoint{3.662121in}{1.802721in}}%
\pgfpathlineto{\pgfqpoint{3.615152in}{1.823962in}}%
\pgfpathlineto{\pgfqpoint{3.568182in}{1.845204in}}%
\pgfpathlineto{\pgfqpoint{3.521212in}{1.866445in}}%
\pgfpathlineto{\pgfqpoint{3.474242in}{1.887687in}}%
\pgfpathlineto{\pgfqpoint{3.427273in}{1.908928in}}%
\pgfpathlineto{\pgfqpoint{3.380303in}{1.930170in}}%
\pgfpathlineto{\pgfqpoint{3.333333in}{1.951411in}}%
\pgfpathlineto{\pgfqpoint{3.286364in}{1.972652in}}%
\pgfpathlineto{\pgfqpoint{3.239394in}{1.993894in}}%
\pgfpathlineto{\pgfqpoint{3.192424in}{2.015135in}}%
\pgfpathlineto{\pgfqpoint{3.145455in}{2.036377in}}%
\pgfpathlineto{\pgfqpoint{3.098485in}{2.057618in}}%
\pgfpathlineto{\pgfqpoint{3.051515in}{2.078860in}}%
\pgfpathlineto{\pgfqpoint{3.004545in}{2.100101in}}%
\pgfpathlineto{\pgfqpoint{2.957576in}{2.121342in}}%
\pgfpathlineto{\pgfqpoint{2.910606in}{2.142584in}}%
\pgfpathlineto{\pgfqpoint{2.863636in}{2.163825in}}%
\pgfpathlineto{\pgfqpoint{2.816667in}{2.185067in}}%
\pgfpathlineto{\pgfqpoint{2.769697in}{2.206308in}}%
\pgfpathlineto{\pgfqpoint{2.722727in}{2.227550in}}%
\pgfpathlineto{\pgfqpoint{2.675758in}{2.248791in}}%
\pgfpathlineto{\pgfqpoint{2.628788in}{2.273453in}}%
\pgfpathlineto{\pgfqpoint{2.581818in}{2.301185in}}%
\pgfpathlineto{\pgfqpoint{2.534848in}{2.328916in}}%
\pgfpathlineto{\pgfqpoint{2.487879in}{2.356648in}}%
\pgfpathlineto{\pgfqpoint{2.440909in}{2.384380in}}%
\pgfpathlineto{\pgfqpoint{2.393939in}{2.412112in}}%
\pgfpathlineto{\pgfqpoint{2.346970in}{2.439844in}}%
\pgfpathlineto{\pgfqpoint{2.300000in}{2.467576in}}%
\pgfpathlineto{\pgfqpoint{2.253030in}{2.495308in}}%
\pgfpathlineto{\pgfqpoint{2.206061in}{2.523039in}}%
\pgfpathlineto{\pgfqpoint{2.159091in}{2.550771in}}%
\pgfpathlineto{\pgfqpoint{2.112121in}{2.578503in}}%
\pgfpathlineto{\pgfqpoint{2.065152in}{2.606235in}}%
\pgfpathlineto{\pgfqpoint{2.018182in}{2.633967in}}%
\pgfpathlineto{\pgfqpoint{1.971212in}{2.661699in}}%
\pgfpathlineto{\pgfqpoint{1.924242in}{2.689431in}}%
\pgfpathlineto{\pgfqpoint{1.877273in}{2.717163in}}%
\pgfpathlineto{\pgfqpoint{1.830303in}{2.744894in}}%
\pgfpathlineto{\pgfqpoint{1.783333in}{2.772626in}}%
\pgfpathlineto{\pgfqpoint{1.736364in}{2.800358in}}%
\pgfpathlineto{\pgfqpoint{1.689394in}{2.828090in}}%
\pgfpathlineto{\pgfqpoint{1.642424in}{2.855822in}}%
\pgfpathlineto{\pgfqpoint{1.595455in}{2.883554in}}%
\pgfpathlineto{\pgfqpoint{1.548485in}{2.911286in}}%
\pgfpathlineto{\pgfqpoint{1.501515in}{2.939017in}}%
\pgfpathlineto{\pgfqpoint{1.454545in}{2.966749in}}%
\pgfpathlineto{\pgfqpoint{1.407576in}{2.994481in}}%
\pgfpathlineto{\pgfqpoint{1.360606in}{3.022213in}}%
\pgfpathlineto{\pgfqpoint{1.313636in}{3.049945in}}%
\pgfpathlineto{\pgfqpoint{1.266667in}{3.077677in}}%
\pgfpathlineto{\pgfqpoint{1.219697in}{3.105409in}}%
\pgfpathlineto{\pgfqpoint{1.172727in}{3.133140in}}%
\pgfpathlineto{\pgfqpoint{1.125758in}{3.160872in}}%
\pgfpathlineto{\pgfqpoint{1.078788in}{3.188604in}}%
\pgfpathlineto{\pgfqpoint{1.031818in}{3.216336in}}%
\pgfpathlineto{\pgfqpoint{0.984848in}{3.244068in}}%
\pgfpathlineto{\pgfqpoint{0.937879in}{3.271800in}}%
\pgfpathlineto{\pgfqpoint{0.890909in}{3.299532in}}%
\pgfpathlineto{\pgfqpoint{0.843939in}{3.327264in}}%
\pgfpathlineto{\pgfqpoint{0.796970in}{3.354995in}}%
\pgfpathlineto{\pgfqpoint{0.750000in}{3.382727in}}%
\pgfpathclose%
\pgfusepath{fill}%
\end{pgfscope}%
\begin{pgfscope}%
\pgfsetbuttcap%
\pgfsetroundjoin%
\definecolor{currentfill}{rgb}{0.000000,0.000000,0.000000}%
\pgfsetfillcolor{currentfill}%
\pgfsetlinewidth{0.803000pt}%
\definecolor{currentstroke}{rgb}{0.000000,0.000000,0.000000}%
\pgfsetstrokecolor{currentstroke}%
\pgfsetdash{}{0pt}%
\pgfsys@defobject{currentmarker}{\pgfqpoint{0.000000in}{-0.048611in}}{\pgfqpoint{0.000000in}{0.000000in}}{%
\pgfpathmoveto{\pgfqpoint{0.000000in}{0.000000in}}%
\pgfpathlineto{\pgfqpoint{0.000000in}{-0.048611in}}%
\pgfusepath{stroke,fill}%
}%
\begin{pgfscope}%
\pgfsys@transformshift{0.967699in}{0.500000in}%
\pgfsys@useobject{currentmarker}{}%
\end{pgfscope}%
\end{pgfscope}%
\begin{pgfscope}%
\definecolor{textcolor}{rgb}{0.000000,0.000000,0.000000}%
\pgfsetstrokecolor{textcolor}%
\pgfsetfillcolor{textcolor}%
\pgftext[x=0.967699in,y=0.402778in,,top]{\color{textcolor}\sffamily\fontsize{10.000000}{12.000000}\selectfont 10}%
\end{pgfscope}%
\begin{pgfscope}%
\pgfsetbuttcap%
\pgfsetroundjoin%
\definecolor{currentfill}{rgb}{0.000000,0.000000,0.000000}%
\pgfsetfillcolor{currentfill}%
\pgfsetlinewidth{0.803000pt}%
\definecolor{currentstroke}{rgb}{0.000000,0.000000,0.000000}%
\pgfsetstrokecolor{currentstroke}%
\pgfsetdash{}{0pt}%
\pgfsys@defobject{currentmarker}{\pgfqpoint{0.000000in}{-0.048611in}}{\pgfqpoint{0.000000in}{0.000000in}}{%
\pgfpathmoveto{\pgfqpoint{0.000000in}{0.000000in}}%
\pgfpathlineto{\pgfqpoint{0.000000in}{-0.048611in}}%
\pgfusepath{stroke,fill}%
}%
\begin{pgfscope}%
\pgfsys@transformshift{1.810619in}{0.500000in}%
\pgfsys@useobject{currentmarker}{}%
\end{pgfscope}%
\end{pgfscope}%
\begin{pgfscope}%
\definecolor{textcolor}{rgb}{0.000000,0.000000,0.000000}%
\pgfsetstrokecolor{textcolor}%
\pgfsetfillcolor{textcolor}%
\pgftext[x=1.810619in,y=0.402778in,,top]{\color{textcolor}\sffamily\fontsize{10.000000}{12.000000}\selectfont 12}%
\end{pgfscope}%
\begin{pgfscope}%
\pgfsetbuttcap%
\pgfsetroundjoin%
\definecolor{currentfill}{rgb}{0.000000,0.000000,0.000000}%
\pgfsetfillcolor{currentfill}%
\pgfsetlinewidth{0.803000pt}%
\definecolor{currentstroke}{rgb}{0.000000,0.000000,0.000000}%
\pgfsetstrokecolor{currentstroke}%
\pgfsetdash{}{0pt}%
\pgfsys@defobject{currentmarker}{\pgfqpoint{0.000000in}{-0.048611in}}{\pgfqpoint{0.000000in}{0.000000in}}{%
\pgfpathmoveto{\pgfqpoint{0.000000in}{0.000000in}}%
\pgfpathlineto{\pgfqpoint{0.000000in}{-0.048611in}}%
\pgfusepath{stroke,fill}%
}%
\begin{pgfscope}%
\pgfsys@transformshift{2.653540in}{0.500000in}%
\pgfsys@useobject{currentmarker}{}%
\end{pgfscope}%
\end{pgfscope}%
\begin{pgfscope}%
\definecolor{textcolor}{rgb}{0.000000,0.000000,0.000000}%
\pgfsetstrokecolor{textcolor}%
\pgfsetfillcolor{textcolor}%
\pgftext[x=2.653540in,y=0.402778in,,top]{\color{textcolor}\sffamily\fontsize{10.000000}{12.000000}\selectfont 14}%
\end{pgfscope}%
\begin{pgfscope}%
\pgfsetbuttcap%
\pgfsetroundjoin%
\definecolor{currentfill}{rgb}{0.000000,0.000000,0.000000}%
\pgfsetfillcolor{currentfill}%
\pgfsetlinewidth{0.803000pt}%
\definecolor{currentstroke}{rgb}{0.000000,0.000000,0.000000}%
\pgfsetstrokecolor{currentstroke}%
\pgfsetdash{}{0pt}%
\pgfsys@defobject{currentmarker}{\pgfqpoint{0.000000in}{-0.048611in}}{\pgfqpoint{0.000000in}{0.000000in}}{%
\pgfpathmoveto{\pgfqpoint{0.000000in}{0.000000in}}%
\pgfpathlineto{\pgfqpoint{0.000000in}{-0.048611in}}%
\pgfusepath{stroke,fill}%
}%
\begin{pgfscope}%
\pgfsys@transformshift{3.496460in}{0.500000in}%
\pgfsys@useobject{currentmarker}{}%
\end{pgfscope}%
\end{pgfscope}%
\begin{pgfscope}%
\definecolor{textcolor}{rgb}{0.000000,0.000000,0.000000}%
\pgfsetstrokecolor{textcolor}%
\pgfsetfillcolor{textcolor}%
\pgftext[x=3.496460in,y=0.402778in,,top]{\color{textcolor}\sffamily\fontsize{10.000000}{12.000000}\selectfont 16}%
\end{pgfscope}%
\begin{pgfscope}%
\pgfsetbuttcap%
\pgfsetroundjoin%
\definecolor{currentfill}{rgb}{0.000000,0.000000,0.000000}%
\pgfsetfillcolor{currentfill}%
\pgfsetlinewidth{0.803000pt}%
\definecolor{currentstroke}{rgb}{0.000000,0.000000,0.000000}%
\pgfsetstrokecolor{currentstroke}%
\pgfsetdash{}{0pt}%
\pgfsys@defobject{currentmarker}{\pgfqpoint{0.000000in}{-0.048611in}}{\pgfqpoint{0.000000in}{0.000000in}}{%
\pgfpathmoveto{\pgfqpoint{0.000000in}{0.000000in}}%
\pgfpathlineto{\pgfqpoint{0.000000in}{-0.048611in}}%
\pgfusepath{stroke,fill}%
}%
\begin{pgfscope}%
\pgfsys@transformshift{4.339381in}{0.500000in}%
\pgfsys@useobject{currentmarker}{}%
\end{pgfscope}%
\end{pgfscope}%
\begin{pgfscope}%
\definecolor{textcolor}{rgb}{0.000000,0.000000,0.000000}%
\pgfsetstrokecolor{textcolor}%
\pgfsetfillcolor{textcolor}%
\pgftext[x=4.339381in,y=0.402778in,,top]{\color{textcolor}\sffamily\fontsize{10.000000}{12.000000}\selectfont 18}%
\end{pgfscope}%
\begin{pgfscope}%
\pgfsetbuttcap%
\pgfsetroundjoin%
\definecolor{currentfill}{rgb}{0.000000,0.000000,0.000000}%
\pgfsetfillcolor{currentfill}%
\pgfsetlinewidth{0.803000pt}%
\definecolor{currentstroke}{rgb}{0.000000,0.000000,0.000000}%
\pgfsetstrokecolor{currentstroke}%
\pgfsetdash{}{0pt}%
\pgfsys@defobject{currentmarker}{\pgfqpoint{0.000000in}{-0.048611in}}{\pgfqpoint{0.000000in}{0.000000in}}{%
\pgfpathmoveto{\pgfqpoint{0.000000in}{0.000000in}}%
\pgfpathlineto{\pgfqpoint{0.000000in}{-0.048611in}}%
\pgfusepath{stroke,fill}%
}%
\begin{pgfscope}%
\pgfsys@transformshift{5.182301in}{0.500000in}%
\pgfsys@useobject{currentmarker}{}%
\end{pgfscope}%
\end{pgfscope}%
\begin{pgfscope}%
\definecolor{textcolor}{rgb}{0.000000,0.000000,0.000000}%
\pgfsetstrokecolor{textcolor}%
\pgfsetfillcolor{textcolor}%
\pgftext[x=5.182301in,y=0.402778in,,top]{\color{textcolor}\sffamily\fontsize{10.000000}{12.000000}\selectfont 20}%
\end{pgfscope}%
\begin{pgfscope}%
\pgfsetbuttcap%
\pgfsetroundjoin%
\definecolor{currentfill}{rgb}{0.000000,0.000000,0.000000}%
\pgfsetfillcolor{currentfill}%
\pgfsetlinewidth{0.803000pt}%
\definecolor{currentstroke}{rgb}{0.000000,0.000000,0.000000}%
\pgfsetstrokecolor{currentstroke}%
\pgfsetdash{}{0pt}%
\pgfsys@defobject{currentmarker}{\pgfqpoint{-0.048611in}{0.000000in}}{\pgfqpoint{0.000000in}{0.000000in}}{%
\pgfpathmoveto{\pgfqpoint{0.000000in}{0.000000in}}%
\pgfpathlineto{\pgfqpoint{-0.048611in}{0.000000in}}%
\pgfusepath{stroke,fill}%
}%
\begin{pgfscope}%
\pgfsys@transformshift{0.750000in}{0.638740in}%
\pgfsys@useobject{currentmarker}{}%
\end{pgfscope}%
\end{pgfscope}%
\begin{pgfscope}%
\definecolor{textcolor}{rgb}{0.000000,0.000000,0.000000}%
\pgfsetstrokecolor{textcolor}%
\pgfsetfillcolor{textcolor}%
\pgftext[x=0.564412in,y=0.585978in,left,base]{\color{textcolor}\sffamily\fontsize{10.000000}{12.000000}\selectfont 0}%
\end{pgfscope}%
\begin{pgfscope}%
\pgfsetbuttcap%
\pgfsetroundjoin%
\definecolor{currentfill}{rgb}{0.000000,0.000000,0.000000}%
\pgfsetfillcolor{currentfill}%
\pgfsetlinewidth{0.803000pt}%
\definecolor{currentstroke}{rgb}{0.000000,0.000000,0.000000}%
\pgfsetstrokecolor{currentstroke}%
\pgfsetdash{}{0pt}%
\pgfsys@defobject{currentmarker}{\pgfqpoint{-0.048611in}{0.000000in}}{\pgfqpoint{0.000000in}{0.000000in}}{%
\pgfpathmoveto{\pgfqpoint{0.000000in}{0.000000in}}%
\pgfpathlineto{\pgfqpoint{-0.048611in}{0.000000in}}%
\pgfusepath{stroke,fill}%
}%
\begin{pgfscope}%
\pgfsys@transformshift{0.750000in}{1.035823in}%
\pgfsys@useobject{currentmarker}{}%
\end{pgfscope}%
\end{pgfscope}%
\begin{pgfscope}%
\definecolor{textcolor}{rgb}{0.000000,0.000000,0.000000}%
\pgfsetstrokecolor{textcolor}%
\pgfsetfillcolor{textcolor}%
\pgftext[x=0.476047in,y=0.983061in,left,base]{\color{textcolor}\sffamily\fontsize{10.000000}{12.000000}\selectfont 25}%
\end{pgfscope}%
\begin{pgfscope}%
\pgfsetbuttcap%
\pgfsetroundjoin%
\definecolor{currentfill}{rgb}{0.000000,0.000000,0.000000}%
\pgfsetfillcolor{currentfill}%
\pgfsetlinewidth{0.803000pt}%
\definecolor{currentstroke}{rgb}{0.000000,0.000000,0.000000}%
\pgfsetstrokecolor{currentstroke}%
\pgfsetdash{}{0pt}%
\pgfsys@defobject{currentmarker}{\pgfqpoint{-0.048611in}{0.000000in}}{\pgfqpoint{0.000000in}{0.000000in}}{%
\pgfpathmoveto{\pgfqpoint{0.000000in}{0.000000in}}%
\pgfpathlineto{\pgfqpoint{-0.048611in}{0.000000in}}%
\pgfusepath{stroke,fill}%
}%
\begin{pgfscope}%
\pgfsys@transformshift{0.750000in}{1.432906in}%
\pgfsys@useobject{currentmarker}{}%
\end{pgfscope}%
\end{pgfscope}%
\begin{pgfscope}%
\definecolor{textcolor}{rgb}{0.000000,0.000000,0.000000}%
\pgfsetstrokecolor{textcolor}%
\pgfsetfillcolor{textcolor}%
\pgftext[x=0.476047in,y=1.380144in,left,base]{\color{textcolor}\sffamily\fontsize{10.000000}{12.000000}\selectfont 50}%
\end{pgfscope}%
\begin{pgfscope}%
\pgfsetbuttcap%
\pgfsetroundjoin%
\definecolor{currentfill}{rgb}{0.000000,0.000000,0.000000}%
\pgfsetfillcolor{currentfill}%
\pgfsetlinewidth{0.803000pt}%
\definecolor{currentstroke}{rgb}{0.000000,0.000000,0.000000}%
\pgfsetstrokecolor{currentstroke}%
\pgfsetdash{}{0pt}%
\pgfsys@defobject{currentmarker}{\pgfqpoint{-0.048611in}{0.000000in}}{\pgfqpoint{0.000000in}{0.000000in}}{%
\pgfpathmoveto{\pgfqpoint{0.000000in}{0.000000in}}%
\pgfpathlineto{\pgfqpoint{-0.048611in}{0.000000in}}%
\pgfusepath{stroke,fill}%
}%
\begin{pgfscope}%
\pgfsys@transformshift{0.750000in}{1.829989in}%
\pgfsys@useobject{currentmarker}{}%
\end{pgfscope}%
\end{pgfscope}%
\begin{pgfscope}%
\definecolor{textcolor}{rgb}{0.000000,0.000000,0.000000}%
\pgfsetstrokecolor{textcolor}%
\pgfsetfillcolor{textcolor}%
\pgftext[x=0.476047in,y=1.777227in,left,base]{\color{textcolor}\sffamily\fontsize{10.000000}{12.000000}\selectfont 75}%
\end{pgfscope}%
\begin{pgfscope}%
\pgfsetbuttcap%
\pgfsetroundjoin%
\definecolor{currentfill}{rgb}{0.000000,0.000000,0.000000}%
\pgfsetfillcolor{currentfill}%
\pgfsetlinewidth{0.803000pt}%
\definecolor{currentstroke}{rgb}{0.000000,0.000000,0.000000}%
\pgfsetstrokecolor{currentstroke}%
\pgfsetdash{}{0pt}%
\pgfsys@defobject{currentmarker}{\pgfqpoint{-0.048611in}{0.000000in}}{\pgfqpoint{0.000000in}{0.000000in}}{%
\pgfpathmoveto{\pgfqpoint{0.000000in}{0.000000in}}%
\pgfpathlineto{\pgfqpoint{-0.048611in}{0.000000in}}%
\pgfusepath{stroke,fill}%
}%
\begin{pgfscope}%
\pgfsys@transformshift{0.750000in}{2.227072in}%
\pgfsys@useobject{currentmarker}{}%
\end{pgfscope}%
\end{pgfscope}%
\begin{pgfscope}%
\definecolor{textcolor}{rgb}{0.000000,0.000000,0.000000}%
\pgfsetstrokecolor{textcolor}%
\pgfsetfillcolor{textcolor}%
\pgftext[x=0.387682in,y=2.174311in,left,base]{\color{textcolor}\sffamily\fontsize{10.000000}{12.000000}\selectfont 100}%
\end{pgfscope}%
\begin{pgfscope}%
\pgfsetbuttcap%
\pgfsetroundjoin%
\definecolor{currentfill}{rgb}{0.000000,0.000000,0.000000}%
\pgfsetfillcolor{currentfill}%
\pgfsetlinewidth{0.803000pt}%
\definecolor{currentstroke}{rgb}{0.000000,0.000000,0.000000}%
\pgfsetstrokecolor{currentstroke}%
\pgfsetdash{}{0pt}%
\pgfsys@defobject{currentmarker}{\pgfqpoint{-0.048611in}{0.000000in}}{\pgfqpoint{0.000000in}{0.000000in}}{%
\pgfpathmoveto{\pgfqpoint{0.000000in}{0.000000in}}%
\pgfpathlineto{\pgfqpoint{-0.048611in}{0.000000in}}%
\pgfusepath{stroke,fill}%
}%
\begin{pgfscope}%
\pgfsys@transformshift{0.750000in}{2.624155in}%
\pgfsys@useobject{currentmarker}{}%
\end{pgfscope}%
\end{pgfscope}%
\begin{pgfscope}%
\definecolor{textcolor}{rgb}{0.000000,0.000000,0.000000}%
\pgfsetstrokecolor{textcolor}%
\pgfsetfillcolor{textcolor}%
\pgftext[x=0.387682in,y=2.571394in,left,base]{\color{textcolor}\sffamily\fontsize{10.000000}{12.000000}\selectfont 125}%
\end{pgfscope}%
\begin{pgfscope}%
\pgfsetbuttcap%
\pgfsetroundjoin%
\definecolor{currentfill}{rgb}{0.000000,0.000000,0.000000}%
\pgfsetfillcolor{currentfill}%
\pgfsetlinewidth{0.803000pt}%
\definecolor{currentstroke}{rgb}{0.000000,0.000000,0.000000}%
\pgfsetstrokecolor{currentstroke}%
\pgfsetdash{}{0pt}%
\pgfsys@defobject{currentmarker}{\pgfqpoint{-0.048611in}{0.000000in}}{\pgfqpoint{0.000000in}{0.000000in}}{%
\pgfpathmoveto{\pgfqpoint{0.000000in}{0.000000in}}%
\pgfpathlineto{\pgfqpoint{-0.048611in}{0.000000in}}%
\pgfusepath{stroke,fill}%
}%
\begin{pgfscope}%
\pgfsys@transformshift{0.750000in}{3.021238in}%
\pgfsys@useobject{currentmarker}{}%
\end{pgfscope}%
\end{pgfscope}%
\begin{pgfscope}%
\definecolor{textcolor}{rgb}{0.000000,0.000000,0.000000}%
\pgfsetstrokecolor{textcolor}%
\pgfsetfillcolor{textcolor}%
\pgftext[x=0.387682in,y=2.968477in,left,base]{\color{textcolor}\sffamily\fontsize{10.000000}{12.000000}\selectfont 150}%
\end{pgfscope}%
\begin{pgfscope}%
\pgfsetbuttcap%
\pgfsetroundjoin%
\definecolor{currentfill}{rgb}{0.000000,0.000000,0.000000}%
\pgfsetfillcolor{currentfill}%
\pgfsetlinewidth{0.803000pt}%
\definecolor{currentstroke}{rgb}{0.000000,0.000000,0.000000}%
\pgfsetstrokecolor{currentstroke}%
\pgfsetdash{}{0pt}%
\pgfsys@defobject{currentmarker}{\pgfqpoint{-0.048611in}{0.000000in}}{\pgfqpoint{0.000000in}{0.000000in}}{%
\pgfpathmoveto{\pgfqpoint{0.000000in}{0.000000in}}%
\pgfpathlineto{\pgfqpoint{-0.048611in}{0.000000in}}%
\pgfusepath{stroke,fill}%
}%
\begin{pgfscope}%
\pgfsys@transformshift{0.750000in}{3.418321in}%
\pgfsys@useobject{currentmarker}{}%
\end{pgfscope}%
\end{pgfscope}%
\begin{pgfscope}%
\definecolor{textcolor}{rgb}{0.000000,0.000000,0.000000}%
\pgfsetstrokecolor{textcolor}%
\pgfsetfillcolor{textcolor}%
\pgftext[x=0.387682in,y=3.365560in,left,base]{\color{textcolor}\sffamily\fontsize{10.000000}{12.000000}\selectfont 175}%
\end{pgfscope}%
\begin{pgfscope}%
\pgfpathrectangle{\pgfqpoint{0.750000in}{0.500000in}}{\pgfqpoint{4.650000in}{3.020000in}}%
\pgfusepath{clip}%
\pgfsetrectcap%
\pgfsetroundjoin%
\pgfsetlinewidth{2.258437pt}%
\definecolor{currentstroke}{rgb}{0.121569,0.466667,0.705882}%
\pgfsetstrokecolor{currentstroke}%
\pgfsetdash{}{0pt}%
\pgfpathmoveto{\pgfqpoint{0.750000in}{3.155920in}}%
\pgfpathlineto{\pgfqpoint{0.796970in}{3.132213in}}%
\pgfpathlineto{\pgfqpoint{0.843939in}{3.108506in}}%
\pgfpathlineto{\pgfqpoint{0.890909in}{3.084800in}}%
\pgfpathlineto{\pgfqpoint{0.937879in}{3.061093in}}%
\pgfpathlineto{\pgfqpoint{0.984848in}{3.037386in}}%
\pgfpathlineto{\pgfqpoint{1.031818in}{3.013680in}}%
\pgfpathlineto{\pgfqpoint{1.078788in}{2.989973in}}%
\pgfpathlineto{\pgfqpoint{1.125758in}{2.966266in}}%
\pgfpathlineto{\pgfqpoint{1.172727in}{2.942560in}}%
\pgfpathlineto{\pgfqpoint{1.219697in}{2.918853in}}%
\pgfpathlineto{\pgfqpoint{1.266667in}{2.895146in}}%
\pgfpathlineto{\pgfqpoint{1.313636in}{2.871440in}}%
\pgfpathlineto{\pgfqpoint{1.360606in}{2.847733in}}%
\pgfpathlineto{\pgfqpoint{1.407576in}{2.824026in}}%
\pgfpathlineto{\pgfqpoint{1.454545in}{2.800320in}}%
\pgfpathlineto{\pgfqpoint{1.501515in}{2.776613in}}%
\pgfpathlineto{\pgfqpoint{1.548485in}{2.752906in}}%
\pgfpathlineto{\pgfqpoint{1.595455in}{2.729200in}}%
\pgfpathlineto{\pgfqpoint{1.642424in}{2.705493in}}%
\pgfpathlineto{\pgfqpoint{1.689394in}{2.681786in}}%
\pgfpathlineto{\pgfqpoint{1.736364in}{2.658080in}}%
\pgfpathlineto{\pgfqpoint{1.783333in}{2.634373in}}%
\pgfpathlineto{\pgfqpoint{1.830303in}{2.610666in}}%
\pgfpathlineto{\pgfqpoint{1.877273in}{2.586960in}}%
\pgfpathlineto{\pgfqpoint{1.924242in}{2.563253in}}%
\pgfpathlineto{\pgfqpoint{1.971212in}{2.539546in}}%
\pgfpathlineto{\pgfqpoint{2.018182in}{2.515840in}}%
\pgfpathlineto{\pgfqpoint{2.065152in}{2.492133in}}%
\pgfpathlineto{\pgfqpoint{2.112121in}{2.468426in}}%
\pgfpathlineto{\pgfqpoint{2.159091in}{2.444720in}}%
\pgfpathlineto{\pgfqpoint{2.206061in}{2.421013in}}%
\pgfpathlineto{\pgfqpoint{2.253030in}{2.397306in}}%
\pgfpathlineto{\pgfqpoint{2.300000in}{2.373600in}}%
\pgfpathlineto{\pgfqpoint{2.346970in}{2.349893in}}%
\pgfpathlineto{\pgfqpoint{2.393939in}{2.326186in}}%
\pgfpathlineto{\pgfqpoint{2.440909in}{2.302480in}}%
\pgfpathlineto{\pgfqpoint{2.487879in}{2.278773in}}%
\pgfpathlineto{\pgfqpoint{2.534848in}{2.255067in}}%
\pgfpathlineto{\pgfqpoint{2.581818in}{2.231360in}}%
\pgfpathlineto{\pgfqpoint{2.628788in}{2.207653in}}%
\pgfpathlineto{\pgfqpoint{2.675758in}{2.183947in}}%
\pgfpathlineto{\pgfqpoint{2.722727in}{2.160240in}}%
\pgfpathlineto{\pgfqpoint{2.769697in}{2.136533in}}%
\pgfpathlineto{\pgfqpoint{2.816667in}{2.112827in}}%
\pgfpathlineto{\pgfqpoint{2.863636in}{2.089120in}}%
\pgfpathlineto{\pgfqpoint{2.910606in}{2.065413in}}%
\pgfpathlineto{\pgfqpoint{2.957576in}{2.041707in}}%
\pgfpathlineto{\pgfqpoint{3.004545in}{2.018000in}}%
\pgfpathlineto{\pgfqpoint{3.051515in}{1.994293in}}%
\pgfpathlineto{\pgfqpoint{3.098485in}{1.970587in}}%
\pgfpathlineto{\pgfqpoint{3.145455in}{1.946880in}}%
\pgfpathlineto{\pgfqpoint{3.192424in}{1.923173in}}%
\pgfpathlineto{\pgfqpoint{3.239394in}{1.899467in}}%
\pgfpathlineto{\pgfqpoint{3.286364in}{1.875760in}}%
\pgfpathlineto{\pgfqpoint{3.333333in}{1.852053in}}%
\pgfpathlineto{\pgfqpoint{3.380303in}{1.828347in}}%
\pgfpathlineto{\pgfqpoint{3.427273in}{1.804640in}}%
\pgfpathlineto{\pgfqpoint{3.474242in}{1.780933in}}%
\pgfpathlineto{\pgfqpoint{3.521212in}{1.757227in}}%
\pgfpathlineto{\pgfqpoint{3.568182in}{1.733520in}}%
\pgfpathlineto{\pgfqpoint{3.615152in}{1.709813in}}%
\pgfpathlineto{\pgfqpoint{3.662121in}{1.686107in}}%
\pgfpathlineto{\pgfqpoint{3.709091in}{1.662400in}}%
\pgfpathlineto{\pgfqpoint{3.756061in}{1.638693in}}%
\pgfpathlineto{\pgfqpoint{3.803030in}{1.614987in}}%
\pgfpathlineto{\pgfqpoint{3.850000in}{1.591280in}}%
\pgfpathlineto{\pgfqpoint{3.896970in}{1.567573in}}%
\pgfpathlineto{\pgfqpoint{3.943939in}{1.543867in}}%
\pgfpathlineto{\pgfqpoint{3.990909in}{1.520160in}}%
\pgfpathlineto{\pgfqpoint{4.037879in}{1.496453in}}%
\pgfpathlineto{\pgfqpoint{4.084848in}{1.472747in}}%
\pgfpathlineto{\pgfqpoint{4.131818in}{1.449040in}}%
\pgfpathlineto{\pgfqpoint{4.178788in}{1.425333in}}%
\pgfpathlineto{\pgfqpoint{4.225758in}{1.401627in}}%
\pgfpathlineto{\pgfqpoint{4.272727in}{1.377920in}}%
\pgfpathlineto{\pgfqpoint{4.319697in}{1.354213in}}%
\pgfpathlineto{\pgfqpoint{4.366667in}{1.330507in}}%
\pgfpathlineto{\pgfqpoint{4.413636in}{1.306800in}}%
\pgfpathlineto{\pgfqpoint{4.460606in}{1.283093in}}%
\pgfpathlineto{\pgfqpoint{4.507576in}{1.259387in}}%
\pgfpathlineto{\pgfqpoint{4.554545in}{1.235680in}}%
\pgfpathlineto{\pgfqpoint{4.601515in}{1.211973in}}%
\pgfpathlineto{\pgfqpoint{4.648485in}{1.188267in}}%
\pgfpathlineto{\pgfqpoint{4.695455in}{1.164560in}}%
\pgfpathlineto{\pgfqpoint{4.742424in}{1.140853in}}%
\pgfpathlineto{\pgfqpoint{4.789394in}{1.117147in}}%
\pgfpathlineto{\pgfqpoint{4.836364in}{1.093440in}}%
\pgfpathlineto{\pgfqpoint{4.883333in}{1.069733in}}%
\pgfpathlineto{\pgfqpoint{4.930303in}{1.046027in}}%
\pgfpathlineto{\pgfqpoint{4.977273in}{1.022320in}}%
\pgfpathlineto{\pgfqpoint{5.024242in}{0.998613in}}%
\pgfpathlineto{\pgfqpoint{5.071212in}{0.974907in}}%
\pgfpathlineto{\pgfqpoint{5.118182in}{0.951200in}}%
\pgfpathlineto{\pgfqpoint{5.165152in}{0.927493in}}%
\pgfpathlineto{\pgfqpoint{5.212121in}{0.903787in}}%
\pgfpathlineto{\pgfqpoint{5.259091in}{0.880080in}}%
\pgfpathlineto{\pgfqpoint{5.306061in}{0.856373in}}%
\pgfpathlineto{\pgfqpoint{5.353030in}{0.832667in}}%
\pgfpathlineto{\pgfqpoint{5.400000in}{0.808960in}}%
\pgfusepath{stroke}%
\end{pgfscope}%
\begin{pgfscope}%
\pgfsetrectcap%
\pgfsetmiterjoin%
\pgfsetlinewidth{0.803000pt}%
\definecolor{currentstroke}{rgb}{0.000000,0.000000,0.000000}%
\pgfsetstrokecolor{currentstroke}%
\pgfsetdash{}{0pt}%
\pgfpathmoveto{\pgfqpoint{0.750000in}{0.500000in}}%
\pgfpathlineto{\pgfqpoint{0.750000in}{3.520000in}}%
\pgfusepath{stroke}%
\end{pgfscope}%
\begin{pgfscope}%
\pgfsetrectcap%
\pgfsetmiterjoin%
\pgfsetlinewidth{0.803000pt}%
\definecolor{currentstroke}{rgb}{0.000000,0.000000,0.000000}%
\pgfsetstrokecolor{currentstroke}%
\pgfsetdash{}{0pt}%
\pgfpathmoveto{\pgfqpoint{5.400000in}{0.500000in}}%
\pgfpathlineto{\pgfqpoint{5.400000in}{3.520000in}}%
\pgfusepath{stroke}%
\end{pgfscope}%
\begin{pgfscope}%
\pgfsetrectcap%
\pgfsetmiterjoin%
\pgfsetlinewidth{0.803000pt}%
\definecolor{currentstroke}{rgb}{0.000000,0.000000,0.000000}%
\pgfsetstrokecolor{currentstroke}%
\pgfsetdash{}{0pt}%
\pgfpathmoveto{\pgfqpoint{0.750000in}{0.500000in}}%
\pgfpathlineto{\pgfqpoint{5.400000in}{0.500000in}}%
\pgfusepath{stroke}%
\end{pgfscope}%
\begin{pgfscope}%
\pgfsetrectcap%
\pgfsetmiterjoin%
\pgfsetlinewidth{0.803000pt}%
\definecolor{currentstroke}{rgb}{0.000000,0.000000,0.000000}%
\pgfsetstrokecolor{currentstroke}%
\pgfsetdash{}{0pt}%
\pgfpathmoveto{\pgfqpoint{0.750000in}{3.520000in}}%
\pgfpathlineto{\pgfqpoint{5.400000in}{3.520000in}}%
\pgfusepath{stroke}%
\end{pgfscope}%
\end{pgfpicture}%
\makeatother%
\endgroup%
}\\La estimación para x=8 es: 178.34703196347   \end{solution}
\end{parts}




\question[2] Una bandeja de doce pimientos de piquillo contiene tres que pican, pero todos tienen el mismo aspecto. Se comen dos pimientos al azar. Cacula:
%\noaddpoints % to omit double points count

\begin{parts}
\part La probabilidad de no coger ninguno que pique
\part  La probabilidad de coger exactamente uno que pique
\part La probabilidad de coger al menos uno que pique
\part La probabilidad de que al coger un tercero, éste pique sabiendo que antes he cogido dos que no picaban 

\end{parts}
\begin{solution}
$[6/11, 47/132, 5/11, 3/10]$
\end{solution}

\addpoints


\question Luis juega a los dardos, y en el 65\% de sus lanzamientos obtiene una puntuación superior a 7 puntos. Sabiendo que en una ronda tira tres dardos, halla la probabilidad de que en una ronda:
%\noaddpoints % to omit double points count

\begin{parts}
\part[1]  Obtenga más de 7 puntos en todos los lanzamientos  \begin{solution}  $0.2746$  \end{solution}
\part[1]  No los obtenga en ninguno  \begin{solution}  $0.04287$  \end{solution}
\part[1]  Si su primer dardo no se clava en la diana, supere los 7 puntos en al menos,
una ocasión  \begin{solution}  $0.8775$  \end{solution}


\end{parts}
\addpoints

\question Una dado se lanza 1000 veces. Calcula la probabilidad de que haya salido número impar: 
%\noaddpoints % to omit double points count
\begin{parts}
\part[1] Al menos 501 veces
\begin{solution} $[0.0316227766016838, 0.5126, 0.4874]$
\\$(500.0, 15.8114, 0.4874, 0.4874)$
\end{solution}

\part[1] Entre 490 y 510 veces, ambos incluidos
\begin{solution}$[0.66407830863536, 0.7467, 0.493362426757813]$\\
$(500.0, 15.8114, 0.4934, 0.4933)$
\end{solution}

\end{parts}
\addpoints


\end{questions}

\newgeometry{left=1 cm,bottom=2cm}
\begin{landscape}
\begin{table}[]
\Large
\centering

\caption{Extracto de tabla de probabilidades de la \textbf{normal estándar $Z(0,1)$}}
\label{my-label}

\begin{tabular}{l|llllllllll}
z   & 0       & 0,01    & 0,02    & 0,03    & 0,04    & 0,05    & 0,06    & 0,07    & 0,08    & 0,09    \\
\hline
0   & 0,5     & 0,50399 & 0,50798 & 0,51197 & 0,51595 & 0,51994 & 0,52392 & 0,5279  & 0,53188 & 0,53586 \\
0,1 & 0,53983 & 0,5438  & 0,54776 & 0,55172 & 0,55567 & 0,55962 & 0,56356 & 0,56749 & 0,57142 & 0,57535 \\
0,2 & 0,57926 & 0,58317 & 0,58706 & 0,59095 & 0,59483 & 0,59871 & 0,60257 & 0,60642 & 0,61026 & 0,61409 \\
0,3 & 0,61791 & 0,62172 & 0,62552 & 0,6293  & 0,63307 & 0,63683 & 0,64058 & 0,64431 & 0,64803 & 0,65173 \\
0,4 & 0,65542 & 0,6591  & 0,66276 & 0,6664  & 0,67003 & 0,67364 & 0,67724 & 0,68082 & 0,68439 & 0,68793 \\
0,5 & 0,69146 & 0,69497 & 0,69847 & 0,70194 & 0,7054  & 0,70884 & 0,71226 & 0,71566 & 0,71904 & 0,7224  \\
0,6 & 0,72575 & 0,72907 & 0,73237 & 0,73565 & 0,73891 & 0,74215 & 0,74537 & 0,74857 & 0,75175 & 0,7549  \\
0,7 & 0,75804 & 0,76115 & 0,76424 & 0,7673  & 0,77035 & 0,77337 & 0,77637 & 0,77935 & 0,7823  & 0,78524 \\
0,8 & 0,78814 & 0,79103 & 0,79389 & 0,79673 & 0,79955 & 0,80234 & 0,80511 & 0,80785 & 0,81057 & 0,81327 \\
0,9 & 0,81594 & 0,81859 & 0,82121 & 0,82381 & 0,82639 & 0,82894 & 0,83147 & 0,83398 & 0,83646 & 0,83891 \\
1   & 0,84134 & 0,84375 & 0,84614 & 0,84849 & 0,85083 & 0,85314 & 0,85543 & 0,85769 & 0,85993 & 0,86214 \\
%1,1 & 0,86433 & 0,8665  & 0,86864 & 0,87076 & 0,87286 & 0,87493 & 0,87698 & 0,879   & 0,881   & 0,88298 \\
%1,2 & 0,88493 & 0,88686 & 0,88877 & 0,89065 & 0,89251 & 0,89435 & 0,89617 & 0,89796 & 0,89973 & 0,90147 \\
%1,3 & 0,9032  & 0,9049  & 0,90658 & 0,90824 & 0,90988 & 0,91149 & 0,91309 & 0,91466 & 0,91621 & 0,91774 \\
%1,4 & 0,91924 & 0,92073 & 0,9222  & 0,92364 & 0,92507 & 0,92647 & 0,92785 & 0,92922 & 0,93056 & 0,93189 \\
%1,5 & 0,93319 & 0,93448 & 0,93574 & 0,93699 & 0,93822 & 0,93943 & 0,94062 & 0,94179 & 0,94295 & 0,94408 \\
%1,6 & 0,9452  & 0,9463  & 0,94738 & 0,94845 & 0,9495  & 0,95053 & 0,95154 & 0,95254 & 0,95352 & 0,95449 \\
%1,7 & 0,95543 & 0,95637 & 0,95728 & 0,95818 & 0,95907 & 0,95994 & 0,9608  & 0,96164 & 0,96246 & 0,96327 \\
%1,8 & 0,96407 & 0,96485 & 0,96562 & 0,96638 & 0,96712 & 0,96784 & 0,96856 & 0,96926 & 0,96995 & 0,97062 \\
%1,9 & 0,97128 & 0,97193 & 0,97257 & 0,9732  & 0,97381 & 0,97441 & 0,975   & 0,97558 & 0,97615 & 0,9767  \\
%2   & 0,97725 & 0,97778 & 0,97831 & 0,97882 & 0,97932 & 0,97982 & 0,9803  & 0,98077 & 0,98124 & 0,98169 \\
%2,1 & 0,98214 & 0,98257 & 0,983   & 0,98341 & 0,98382 & 0,98422 & 0,98461 & 0,985   & 0,98537 & 0,98574 \\
%2,2 & 0,9861  & 0,98645 & 0,98679 & 0,98713 & 0,98745 & 0,98778 & 0,98809 & 0,9884  & 0,9887  & 0,98899 \\
%2,3 & 0,98928 & 0,98956 & 0,98983 & 0,9901  & 0,99036 & 0,99061 & 0,99086 & 0,99111 & 0,99134 & 0,99158 \\
%2,4 & 0,9918  & 0,99202 & 0,99224 & 0,99245 & 0,99266 & 0,99286 & 0,99305 & 0,99324 & 0,99343 & 0,99361 \\
%2,5 & 0,99379 & 0,99396 & 0,99413 & 0,9943  & 0,99446 & 0,99461 & 0,99477 & 0,99492 & 0,99506 & 0,9952  \\
%2,6 & 0,99534 & 0,99547 & 0,9956  & 0,99573 & 0,99585 & 0,99598 & 0,99609 & 0,99621 & 0,99632 & 0,99643 \\
%2,7 & 0,99653 & 0,99664 & 0,99674 & 0,99683 & 0,99693 & 0,99702 & 0,99711 & 0,9972  & 0,99728 & 0,99736 \\
%2,8 & 0,99744 & 0,99752 & 0,9976  & 0,99767 & 0,99774 & 0,99781 & 0,99788 & 0,99795 & 0,99801 & 0,99807 \\
%2,9 & 0,99813 & 0,99819 & 0,99825 & 0,99831 & 0,99836 & 0,99841 & 0,99846 & 0,99851 & 0,99856 & 0,99861 \\
%3   & 0,99865 & 0,99869 & 0,99874 & 0,99878 & 0,99882 & 0,99886 & 0,99889 & 0,99893 & 0,99896 & 0,999   \\
%3,1 & 0,99903 & 0,99906 & 0,9991  & 0,99913 & 0,99916 & 0,99918 & 0,99921 & 0,99924 & 0,99926 & 0,99929 \\
%3,2 & 0,99931 & 0,99934 & 0,99936 & 0,99938 & 0,9994  & 0,99942 & 0,99944 & 0,99946 & 0,99948 & 0,9995  \\
%3,3 & 0,99952 & 0,99953 & 0,99955 & 0,99957 & 0,99958 & 0,9996  & 0,99961 & 0,99962 & 0,99964 & 0,99965 \\
%3,4 & 0,99966 & 0,99968 & 0,99969 & 0,9997  & 0,99971 & 0,99972 & 0,99973 & 0,99974 & 0,99975 & 0,99976 \\
%3,5 & 0,99977 & 0,99978 & 0,99978 & 0,99979 & 0,9998  & 0,99981 & 0,99981 & 0,99982 & 0,99983 & 0,99983 \\
%3,6 & 0,99984 & 0,99985 & 0,99985 & 0,99986 & 0,99986 & 0,99987 & 0,99987 & 0,99988 & 0,99988 & 0,99989 \\
%3,7 & 0,99989 & 0,9999  & 0,9999  & 0,9999  & 0,99991 & 0,99991 & 0,99992 & 0,99992 & 0,99992 & 0,99992 \\
%3,8 & 0,99993 & 0,99993 & 0,99993 & 0,99994 & 0,99994 & 0,99994 & 0,99994 & 0,99995 & 0,99995 & 0,99995 \\
%3,9 & 0,99995 & 0,99995 & 0,99996 & 0,99996 & 0,99996 & 0,99996 & 0,99996 & 0,99996 & 0,99997 & 0,99997 \\
%4   & 0,99997 & 0,99997 & 0,99997 & 0,99997 & 0,99997 & 0,99997 & 0,99998 & 0,99998 & 0,99998 & 0,99998
\end{tabular}
\end{table}
\end{landscape}
\restoregeometry

\end{document}
\grid
